% Target: rmp-workbench-iii-g-injective-pull
% Formula: A virtual g-string pulls through a G-injective PEPS tensor.
% Ink: Canonical public kernel PEPS star with opposite g-string corners.
% Boundary: Five PEPS legs and two g-string ends remain open on each side.
% Source: author source ImagesReview Section III/ImagesReview Section III.tex
%   lines 468-497; archival output Ginjpullt.pdf
% Capabilities: kernel, pulling-through, open-string
\[
  % Author panels (lines 468-497): the pulling-through equality for the
  % G-injective PEPS tensor A.  A virtual g-string crosses the west and
  % north legs, a group-action box at each crossing; pulled through, it
  % crosses the east and south legs instead.  Both physical legs stay short.
  \tenkzkernel
  \tndeclare{species}{g}{hue=source:red}
  \begin{tenkzeq}[size=m]
    \begin{tenkz}[rows={wire,wire,wire,wire,wire}, cols=5, bonds=none]
      \tn[at=(3,3), name=A, ports={0:virtual, 90:virtual, 180:virtual,
          270:virtual, 45:physical}]{}
      \tnwire[kind=index]{A.180}{(3,1)}
      \tnwire[kind=index]{A.0}{(3,5)}
      \tnwire[kind=index]{A.270}{(5,3)}
      \tnwire[kind=index]{A.90}{(1,3)}
      % the g-string's two waypoints become invisible junction atoms, chained
      % by wires named after the original single wire (via= retired, #6897)
      \tn[at=(3,2), skin=none, name=jg1]{}
      \tn[at=(2,3), skin=none, name=jg2]{}
      \tnwire[kind=string, species=g, name=g]{(4,2)}{jg1}
      \tnwire[kind=string, species=g, name=g2,
        cross={under at crossing of g and g2}]{jg1}{jg2}
      \tnwire[kind=string, species=g, name=g3,
        cross={under at crossing of g2 and g3}]{jg2}{(2,5)}
      \tn[skin=box, at=on g2 0.04]{}
      \tn[skin=box, at=on g3 0.03]{}
      \tnmark[form=label]{(2,2)}{$g$}
    \end{tenkz}
    =
    \begin{tenkz}[rows={wire,wire,wire,wire,wire}, cols=5, bonds=none]
      \tn[at=(3,3), name=A, ports={0:virtual, 90:virtual, 180:virtual,
          270:virtual, 45:physical}]{}
      \tnwire[kind=index]{A.180}{(3,1)}
      \tnwire[kind=index]{A.0}{(3,5)}
      \tnwire[kind=index]{A.270}{(5,3)}
      \tnwire[kind=index]{A.90}{(1,3)}
      \tn[at=(3,4), skin=none, name=jg1]{}
      \tn[at=(4,3), skin=none, name=jg2]{}
      \tnwire[kind=string, species=g, name=g]{(2,5)}{jg1}
      \tnwire[kind=string, species=g, name=g2]{jg1}{jg2}
      \tnwire[kind=string, species=g, name=g3,
        cross={under at crossing of g2 and g3}]{jg2}{(4,2)}
      \tn[skin=box, at=on g2 0.012]{}
      \tn[skin=box, at=on g3 0.02]{}
      \tnmark[form=label]{(4,4)}{$g$}
    \end{tenkz}
  \end{tenkzeq}
\]
